% GENERATED FILE. Edit canonical lessons, then run npm run generate:books.
% canonical-source-hash: fnv1a64:9f4959d2d81b85c8
% canonical-lessons: SA-C232-vartate, SA-C232-jayate, SA-C232-mriyate, SA-C232-nasyati, SA-C232-kupyati

\chapter{Is present, Is born, Dies, Is lost, Gets angry}
\label{ch:sa-is-present-is-born-dies-is-lost-gets-angry}
\hlchaptermodality{232}

\begin{chapteropening}
\textbf{By the end of this chapter:} \emph{I can say is present, is born, dies, is lost and gets angry.}

Complete the last lesson of chapter 232: I can say is present, is born, dies, is lost and gets angry.
\end{chapteropening}

\section[\texorpdfstring{\sk{वर्तते} (vartate)}{वर्तते (vartate)}]{\sk{वर्तते} (vartate) — is present, exists}

\label{lesson:SA-C232-vartate}



\begin{quote}
Before the new one: say the Sanskrit for a guest house, then the Sanskrit for a stay abroad, a journey.
\end{quote}

\subsection*{You'll want to know: \sk{वर्तते}}
\textbf{\sk{वर्तते}} — \emph{vartate} — \textquotedblleft{}is present, exists\textquotedblright{}.

\begin{grammarlens}[title={What you've built}]
One more word for this chapter.
\end{grammarlens}

\subsection*{Your turn}
\begin{itemize}
\raggedright
  \item \emph{Say it:} \emph{vartate}
  \item \emph{Say it:} \emph{vartate}, once more
  \item \emph{Say it:} say \emph{vartate}
  \item \emph{Recall:} say the Sanskrit for a guest house, then the Sanskrit for a stay abroad, a journey, then say \emph{vartate} again
\end{itemize}

\begin{tcolorbox}[breakable,colback=teal!4,colframe=teal!35!black,title={Before you move on}]
What does \sk{वर्तते} mean? (\textquotedblleft{}Is present, exists\textquotedblright{}.) Say it once more.
\end{tcolorbox}
\section[\texorpdfstring{\sk{जायते} (jāyate)}{जायते (jāyate)}]{\sk{जायते} (jāyate) — is born, arises}

\label{lesson:SA-C232-jayate}



\begin{quote}
Before the new one: say the Sanskrit for a stay abroad, a journey, then the Sanskrit for is present.
\end{quote}

\subsection*{You'll want to know: \sk{जायते}}
\textbf{\sk{जायते}} — \emph{jāyate} — \textquotedblleft{}is born, arises\textquotedblright{}.

\begin{grammarlens}[title={What you've built}]
One more word for this chapter.
\end{grammarlens}

\subsection*{Your turn}
\begin{itemize}
\raggedright
  \item \emph{Say it:} \emph{jāyate}
  \item \emph{Say it:} \emph{jāyate}, once more
  \item \emph{Say it:} say \emph{jāyate}
  \item \emph{Recall:} say the Sanskrit for a stay abroad, a journey, then the Sanskrit for is present, then say \emph{jāyate} again
\end{itemize}

\begin{tcolorbox}[breakable,colback=teal!4,colframe=teal!35!black,title={Before you move on}]
What does \sk{जायते} mean? (\textquotedblleft{}Is born, arises\textquotedblright{}.) Say it once more.
\end{tcolorbox}
\section[\texorpdfstring{\sk{म्रियते} (mriyate)}{म्रियते (mriyate)}]{\sk{म्रियते} (mriyate) — dies}

\label{lesson:SA-C232-mriyate}



\begin{quote}
Before the new one: say the Sanskrit for is present, then the Sanskrit for is born.
\end{quote}

\subsection*{You'll want to know: \sk{म्रियते}}
\textbf{\sk{म्रियते}} — \emph{mriyate} — \textquotedblleft{}dies\textquotedblright{}.

\begin{grammarlens}[title={What you've built}]
One more word for this chapter.
\end{grammarlens}

\subsection*{Your turn}
\begin{itemize}
\raggedright
  \item \emph{Say it:} \emph{mriyate}
  \item \emph{Say it:} \emph{mriyate}, once more
  \item \emph{Say it:} say \emph{mriyate}
  \item \emph{Recall:} say the Sanskrit for is present, then the Sanskrit for is born, then say \emph{mriyate} again
\end{itemize}

\begin{tcolorbox}[breakable,colback=teal!4,colframe=teal!35!black,title={Before you move on}]
What does \sk{म्रियते} mean? (\textquotedblleft{}Dies\textquotedblright{}.) Say it once more.
\end{tcolorbox}
\section[\texorpdfstring{\sk{नश्यति} (naśyati)}{नश्यति (naśyati)}]{\sk{नश्यति} (naśyati) — is lost, perishes}

\label{lesson:SA-C232-nasyati}



\begin{quote}
Before the new one: say the Sanskrit for is born, then the Sanskrit for dies.
\end{quote}

\subsection*{You'll want to know: \sk{नश्यति}}
\textbf{\sk{नश्यति}} — \emph{naśyati} — \textquotedblleft{}is lost, perishes\textquotedblright{}.

\begin{grammarlens}[title={What you've built}]
One more word for this chapter.
\end{grammarlens}

\subsection*{Your turn}
\begin{itemize}
\raggedright
  \item \emph{Say it:} \emph{naśyati}
  \item \emph{Say it:} \emph{naśyati}, once more
  \item \emph{Say it:} say \emph{naśyati}
  \item \emph{Recall:} say the Sanskrit for is born, then the Sanskrit for dies, then say \emph{naśyati} again
\end{itemize}

\begin{tcolorbox}[breakable,colback=teal!4,colframe=teal!35!black,title={Before you move on}]
What does \sk{नश्यति} mean? (\textquotedblleft{}Is lost, perishes\textquotedblright{}.) Say it once more.
\end{tcolorbox}
\section[\texorpdfstring{\sk{कुप्यति} (kupyati)}{कुप्यति (kupyati)}]{\sk{कुप्यति} (kupyati) — gets angry}

\label{lesson:SA-C232-kupyati}



\begin{quote}
Before the new one: say the Sanskrit for dies, then the Sanskrit for is lost.
\end{quote}

\subsection*{You'll want to know: \sk{कुप्यति}}
\textbf{\sk{कुप्यति}} — \emph{kupyati} — \textquotedblleft{}gets angry\textquotedblright{}.

\begin{grammarlens}[title={What you've built}]
One more word for this chapter.
\end{grammarlens}

\subsection*{Your turn}
\begin{itemize}
\raggedright
  \item \emph{Say it:} \emph{kupyati}
  \item \emph{Say it:} \emph{kupyati}, once more
  \item \emph{Say it:} say \emph{kupyati}
  \item \emph{Recall:} say the Sanskrit for dies, then the Sanskrit for is lost, then say \emph{kupyati} again
\end{itemize}

\begin{tcolorbox}[breakable,colback=teal!4,colframe=teal!35!black,title={Before you move on}]
What does \sk{कुप्यति} mean? (\textquotedblleft{}Gets angry\textquotedblright{}.) Say it once more.
\end{tcolorbox}
